\Needspace{12\baselineskip}
\part{From finite projections to continuation}

\section{The compatible family as an estimate generator}
\label{sec:whole-terminal-consequences}

Part~III changed the estimate problem.  Once a finite derivative and target
norm have been named, equation~\eqref{eq:finite-projection-mixed-norms}
selects one sufficiently deep rectangle and supplies the required bound.
The calculations below return to the quantities that made the terminal threat
visible in Parts~I and~II: the classical continuation coefficient, the
Ladyzhenskaya--Prodi--Serrin scale, vorticity and strain, normalized pressure,
the critical height, and the high-frequency tail.  Each comes from a finite
projection of the same compatible family.

These bounds will exclude Blown and the critical escape scenario.  They do not
yet prove that traces obtained in different finite topologies are one terminal
state, or that the normalized pressure--momentum response survives there.
Those are the remaining Jump and Dead obligations.  The finite readouts will
lead directly to them.

On a proposed finite maximal interval,
Theorem~\ref{thm:whole-terminal} at the rectangle
\((N,M)=(0,2)\) gives
\[
 u\in L^2(I_*;H^3).
\]
The embedding \(H^3(\R^3)\hookrightarrow W^{1,\infty}(\R^3)\) and
Cauchy--Schwarz in time yield
\begin{equation}
 L_*:=\int_0^{T_*}\norm{\nabla u(t)}_\infty\,dt
 \le C T_*^{1/2}\norm{u}_{L^2(I_*;H^3)}
 \le C T_*^{1/2}\mathcal N_{0,2}(u;I_*)^{1/2}<\infty.
 \label{eq:terminal-lipschitz-bound}
\end{equation}
This is exactly the quantity whose divergence was forced by
\eqref{eq:lipschitz-blowup-criterion}, and it is the coefficient in the
high-order energy inequality \eqref{eq:high-order-continuation}.  Substituting
\eqref{eq:terminal-lipschitz-bound} into
\eqref{eq:high-order-gronwall} gives, for each fixed integer \(k\ge3\),
\begin{equation}
 \sup_{0\le t<T_*}\norm{u(t)}_{H^k}^2
 \le \norm{u_0}_{H^k}^2\exp(2C_kL_*)
 \le \norm{u_0}_{H^k}^2
 \exp\!\left(
  2C_kC T_*^{1/2}\mathcal N_{0,2}(u;I_*)^{1/2}
 \right)<\infty.
 \label{eq:terminal-sobolev-bound}
\end{equation}
The constant may grow with \(k\).  The conclusion is one finite bound after
one finite order has been chosen.

This pays the direct continuation coefficient isolated in Part~I.  It is
finite on the complete proposed terminal interval, and
\eqref{eq:terminal-sobolev-bound} propagates the resulting control through
every fixed Sobolev order.

\subsection{The numerical size of the completed family}

Theorem~\ref{thm:whole-terminal} has already supplied every compatible
rectangle.  We can now ask how much of that family must remain visible in the
constants of the final estimates.  Put
\begin{equation}
 \mathfrak B_2[u_0,\nu;T_*]
 :=\mathcal N_{0,2}(u;I_*)
 =\norm u_{L^2(I_*;H^3)}^2
  +\norm{u_t}_{L^2(I_*;H^1)}^2.
 \label{eq:base-producer-value}
\end{equation}
This is the finite \((0,2)\) output of the datum theorem.  The calculation
below does not use it to construct that theorem; it shows that, once the
family exists, every higher numerical readout can be written through this one
datum-generated value and finitely many Sobolev sizes of \(u_0\).

The equation itself opens this base value one step further.  Let
\[
 A:=-\Delta,
 \qquad
 B(u,u):=\PP\nabla\cdot(u\otimes u),
\]
so that the projected momentum equation is
\(u_t+\nu Au+B(u,u)=0\).  For \(0<S<T_*\), put
\[
 X_\nu(S)
 :=\int_0^S\left(
   \norm{u_t(t)}_{H^1}^2
   +\nu^2\norm{Au(t)}_{H^1}^2
  \right)dt,
\]
\[
 Z_0
 :=\norm{A^{1/2}u_0}_{H^1}^2
 =\norm{\nabla u_0}_2^2+\norm{\Delta u_0}_2^2,
 \qquad
 J_1(S)
 :=\int_0^S\norm{B(u(t),u(t))}_{H^1}^2\,dt.
\]
Squaring the projected equation in \(H^1\) leaves one cross term.  The
self-adjointness of \(A\) turns it into an endpoint difference:
\begin{align}
 J_1(S)
 &=X_\nu(S)+2\nu\int_0^S\pair{u_t}{Au}_{H^1}\,dt
 \notag\\
 &=X_\nu(S)
   +\nu\left(
    \norm{A^{1/2}u(S)}_{H^1}^2-Z_0
   \right).
 \label{eq:base-graph-identity}
\end{align}
In particular,
\begin{equation}
 X_\nu(S)\le\nu Z_0+J_1(S).
 \label{eq:base-graph-upper-bound}
\end{equation}
The elementary Fourier inequality
\((1+|\xi|^2)^3\le
4(1+|\xi|^4+|\xi|^6)\), together with the kinetic-energy bound, now gives
\begin{equation}
 \mathfrak B_2[u_0,\nu;T_*]
 \le4T_*\norm{u_0}_2^2
 +\kappa_\nu\left(
   \nu Z_0+\mathfrak J_1[u_0,\nu;T_*]
  \right),
 \qquad
 \kappa_\nu:=\max\{1,4\nu^{-2}\},
 \label{eq:base-producer-graph-bound}
\end{equation}
where
\[
 \mathfrak J_1[u_0,\nu;T_*]
 :=\int_0^{T_*}\norm{\PP\nabla\cdot(u\otimes u)}_{H^1}^2\,dt.
\]
This nonlinear action is evaluated on the unique history generated by the
full datum.  Equation~\eqref{eq:base-producer-graph-bound} exposes its exact
place in the base rectangle.  The signed critical current gives a second
description that is closer to the height mechanism from Part~II.  Define
\begin{equation}
 \mathfrak C_{1/2}[u_0,\nu;T_*]
 :=\sup_{0\le S<T_*}\int_0^S\mathcal P_{1/2}(t)\,dt.
 \label{eq:critical-prefix-value}
\end{equation}
The value at \(S=0\) makes \(\mathfrak C_{1/2}\ge0\).  The whole-terminal
theorem makes it finite; the direct \(L^1\) calculation appears in
\eqref{eq:quantitative-critical-response}.  Thus
\(\mathfrak C_{1/2}\) is a scalar value generated by the same datum and
history, not an additional hypothesis.

The exact critical balance can now be read numerically.  For every
\(0<S<T_*\),
\begin{equation*}
 H(S)+\nu\int_0^S\norm{\Lambda^{3/2}u(t)}_2^2\,dt
 =H(0)+\int_0^S\mathcal P_{1/2}(t)\,dt.
\end{equation*}
Since \(H(S)\ge0\), monotone convergence gives
\begin{equation}
 \int_0^{T_*}\norm{\Lambda^{3/2}u(t)}_2^2\,dt
 \le \mathfrak A_{1/2}[u_0,\nu;T_*]
 :=\frac{H(0)+\mathfrak C_{1/2}[u_0,\nu;T_*]}{\nu},
 \qquad
 H(0)=\frac12\norm{\Lambda^{1/2}u_0}_2^2.
 \label{eq:critical-coefficient-value}
\end{equation}
The base rectangle supplies the same coefficient without passing through the
signed current:
\begin{equation}
 \int_0^{T_*}\norm{\Lambda^{3/2}u(t)}_2^2\,dt
 \le \norm u_{L^2(I_*;H^3)}^2
 \le \mathfrak B_2[u_0,\nu;T_*].
 \label{eq:base-rectangle-critical-coefficient}
\end{equation}
The kinetic-energy payment sharpens this direct inclusion.  Fourier
log-convexity at each time and H\"older's inequality in time give
\begin{align*}
 \int_0^{T_*}\norm{\Lambda^{3/2}u(t)}_2^2\,dt
 &\le
 \left(\int_0^{T_*}\norm{\Lambda u(t)}_2^2\,dt\right)^{3/4}
 \left(\int_0^{T_*}\norm{\Lambda^3u(t)}_2^2\,dt\right)^{1/4}\\
 &\le
 \left(\frac{\norm{u_0}_2^2}{2\nu}\right)^{3/4}
 \mathfrak B_2[u_0,\nu;T_*]^{1/4}.
\end{align*}
This prefix bound is already paid by the same base rectangle.  On every
strict subinterval, the critical balance
\(\mathcal P_{1/2}=H'+2\nu E_{3/2}\) and Cauchy--Schwarz give
\begin{align*}
 \int_0^S|H'(t)|\,dt
 &\le
 \norm{u_t}_{L^2(0,S;H^1)}
 \norm{u}_{L^2(0,S;H^1)}
 \le\frac12\mathfrak B_2[u_0,\nu;T_*],\\
 2\nu\int_0^S E_{3/2}(t)\,dt
 &\le\nu\mathfrak B_2[u_0,\nu;T_*].
\end{align*}
Taking the supremum in \(S\) gives
\begin{equation}
 \mathfrak C_{1/2}[u_0,\nu;T_*]
 \le\norm{\mathcal P_{1/2}}_{L^1(I_*)}
 \le\left(\frac12+\nu\right)
 \mathfrak B_2[u_0,\nu;T_*].
 \label{eq:base-rectangle-critical-prefix}
\end{equation}
The strongest current-free coefficient obtained from these two payments is
\begin{equation}
 \mathfrak A_{1/2}^{\rm prod}[u_0,\nu;T_*]
 :=\min\left\{
 \mathfrak B_2[u_0,\nu;T_*],
 \left(\frac{\norm{u_0}_2^2}{2\nu}\right)^{3/4}
 \mathfrak B_2[u_0,\nu;T_*]^{1/4}
 \right\}.
 \label{eq:producer-critical-coefficient}
\end{equation}
Then
\begin{equation}
 \int_0^{T_*}\norm{\Lambda^{3/2}u(t)}_2^2\,dt
 \le \mathfrak A_{1/2}^{\rm prod}[u_0,\nu;T_*].
 \label{eq:producer-critical-coefficient-bound}
\end{equation}
Thus every spatial row is driven directly by the single \((0,2)\) rectangle;
the signed prefix remains an optional sharper readout, not a parameter needed
by the final estimates.
For \(0<T\le T_*\), the same symbols with \(T\) in the final slot denote
the definitions above with \((0,T)\) in place of \(I_*\).

\begin{proposition}[Datum-only calibrations of the critical coefficient]
\label{prop:critical-coefficient-calibrations}
Let \(C_{\rm crit}\) be the constant in
\eqref{eq:critical-response-estimate}, and put
\[
 q_0:=\norm{\Lambda^{1/2}u_0}_2,
 \qquad
 \vartheta:=\frac{C_{\rm crit}q_0}{\nu}.
\]
If \(\vartheta<1\), then every finite horizon satisfies
\begin{equation}
 \mathfrak C_{1/2}[u_0,\nu;T]
 \le\frac{\vartheta}{1-\vartheta}H(0),
 \qquad
 \mathfrak A_{1/2}[u_0,\nu;T]
 \le\frac{H(0)}{\nu(1-\vartheta)}.
 \label{eq:small-critical-coefficient}
\end{equation}

For arbitrary datum, let \(C_{\rm ens}\) be the constant in
\eqref{eq:enstrophy-obstruction}, set
\[
 E_0:=\norm{u_0}_2^2,
 \qquad
 D_1:=\norm{\Lambda u_0}_2^2,
 \qquad
 \Theta(T):=1-2C_{\rm ens}\nu^{-3}D_1^2T,
\]
and suppose \(\Theta(T)>0\).  Then
\begin{align}
 \sup_{0\le t\le T}\norm{\Lambda u(t)}_2^2
 &\le D_1\Theta(T)^{-1/2},
 \label{eq:local-enstrophy-envelope}\\
 \int_0^T\norm{\Lambda^{3/2}u(t)}_2^2\,dt
 &\le\frac{\sqrt{E_0D_1/2}}
 {\nu\Theta(T)^{1/4}},
 \label{eq:local-critical-action}\\
 \mathfrak C_{1/2}[u_0,\nu;T]
 &\le\left(\frac12+\frac1{\sqrt2}\right)
 \sqrt{E_0D_1}\,\Theta(T)^{-1/4}.
 \label{eq:local-critical-prefix}
\end{align}
\end{proposition}

\begin{proof}
When \(\vartheta<1\), estimate
\eqref{eq:critical-response-estimate} and a first-crossing argument give
\[
 H(S)+\nu(1-\vartheta)
 \int_0^S\norm{\Lambda^{3/2}u(t)}_2^2\,dt
 \le H(0).
\]
Writing the signed prefix as
\[
 \left[
 H(S)-H(0)+\nu(1-\vartheta)
 \int_0^S\norm{\Lambda^{3/2}u}_2^2\,dt
 \right]
 +\nu\vartheta
 \int_0^S\norm{\Lambda^{3/2}u}_2^2\,dt
\]
gives \eqref{eq:small-critical-coefficient}.

For the second statement, put
\(Y=\norm{\Lambda u}_2^2\) and
\(Z=\norm{\Lambda^2u}_2^2\).  Scalar comparison in
\eqref{eq:enstrophy-obstruction} gives
\[
 Y(t)\le D_1
 \left(1-2C_{\rm ens}\nu^{-3}D_1^2t\right)^{-1/2}.
\]
Integrating the same inequality and inserting this comparison yields
\[
 \int_0^T Z(t)\,dt
 \le\frac{D_1}{\nu\Theta(T)^{1/2}}.
\]
Fourier interpolation, Cauchy--Schwarz in time, and the kinetic-energy
identity now give
\[
 \int_0^T\norm{\Lambda^{3/2}u}_2^2\,dt
 \le
 \left(\int_0^T Y\,dt\right)^{1/2}
 \left(\int_0^T Z\,dt\right)^{1/2}
 \le\frac{\sqrt{E_0D_1/2}}{\nu\Theta(T)^{1/4}}.
\]
Finally,
\[
 H(S)\le\frac12\norm{u(S)}_2\norm{\Lambda u(S)}_2
 \le\frac12\sqrt{E_0D_1}\,\Theta(T)^{-1/4}.
\]
Insert the last two estimates into the exact critical balance and take the
supremum over \(S\le T\) to obtain \eqref{eq:local-critical-prefix}.
\end{proof}

The preceding estimate uses the first differentiated row and is strongest near
the local enstrophy scale.  A higher datum row offers a different trade.  It
can be very large, but it enters with a fractional exponent that decreases as
the chosen order rises.  This produces a closed datum condition in which the
spatial order, the horizon, and the viscosity remain visible separately.

\begin{proposition}[A finite-order datum window]
\label{prop:finite-order-datum-window}
Fix an integer \(\ell\ge3\).  Let \(C_\ell^{\rm Lip}\) be a constant for
which the \(H^\ell\) transport estimate gives
\begin{equation}
 \norm{u(t)}_{H^\ell}
 \le \norm{u_0}_{H^\ell}
 \exp\!\left(C_\ell^{\rm Lip}\mathcal L(t)\right),
 \qquad
 \mathcal L(t):=\int_0^t\norm{\nabla u(s)}_\infty\,ds.
 \label{eq:finite-order-lipschitz-persistence}
\end{equation}
Put
\begin{equation}
 \theta_\ell:=\frac{3}{2(\ell-1)},
 \qquad
 \lambda_\ell:=C_\ell^{\rm Lip}\theta_\ell,
 \label{eq:finite-order-window-exponents}
\end{equation}
and choose \(G_\ell\) so that
\begin{equation}
 \norm{\nabla f}_\infty
 \le G_\ell
 \norm{\Lambda f}_2^{1-\theta_\ell}
 \norm f_{H^\ell}^{\theta_\ell}.
 \label{eq:finite-order-fourier-split}
\end{equation}
For \(T>0\), define the closed datum value
\begin{equation}
 \Xi_\ell[u_0,\nu;T]
 :=\lambda_\ell G_\ell
 \norm{u_0}_{H^\ell}^{\theta_\ell}
 T^{(1+\theta_\ell)/2}
 \left(\frac{\norm{u_0}_2^2}{2\nu}\right)^{(1-\theta_\ell)/2}.
 \label{eq:finite-order-window-value}
\end{equation}
If \(\Xi_\ell[u_0,\nu;T]<1\), then the solution continues through \(T\)
and
\begin{equation}
 \mathcal L(T)
 \le \frac1{\lambda_\ell}
 \log\frac1{1-\Xi_\ell[u_0,\nu;T]}.
 \label{eq:finite-order-lipschitz-clock}
\end{equation}

Let \(C_1^{\rm Lip}\) be a constant in the first-order estimate
\begin{equation}
 Y'(t)+2\nu Z(t)
 \le2C_1^{\rm Lip}\norm{\nabla u(t)}_\infty Y(t),
 \qquad
 Y:=\norm{\Lambda u}_2^2,
 \quad Z:=\norm{\Lambda^2u}_2^2.
 \label{eq:lipschitz-enstrophy-row}
\end{equation}
With \(E_0=\norm{u_0}_2^2\) and
\(D_1=\norm{\Lambda u_0}_2^2\), the same datum window gives
\begin{equation}
\begin{aligned}
 \mathfrak C_{1/2}[u_0,\nu;T]
 &\le \widehat{\mathfrak C}_{1/2}^{(\ell)}[u_0,\nu;T]
 \\
 &:=\sqrt{E_0D_1}\,
 \bigl(1-\Xi_\ell[u_0,\nu;T]\bigr)^{
 -C_1^{\rm Lip}/\lambda_\ell}
 -H(0).
\end{aligned}
\label{eq:finite-order-critical-prefix}
\end{equation}
\begin{equation}
\begin{aligned}
 \mathfrak A_{1/2}[u_0,\nu;T]
 &\le \widehat{\mathfrak A}_{1/2}^{(\ell)}[u_0,\nu;T]
 :=\frac{\sqrt{E_0D_1}}{\nu}
 \bigl(1-\Xi_\ell[u_0,\nu;T]\bigr)^{
 -C_1^{\rm Lip}/\lambda_\ell}.
\end{aligned}
\label{eq:finite-order-critical-coefficient}
\end{equation}
Every quantity on the right is calculated directly from
\(u_0,\nu,T,\ell\), and the fixed analytic constants.
\end{proposition}

\begin{proof}
Fourier inversion gives
\[
 \norm{\nabla f}_\infty
 \le C\int_{\R^3}|\xi|\,|\widehat f(\xi)|\,d\xi.
\]
Split the integral at a radius \(R>0\).  Cauchy--Schwarz on the two pieces
gives
\[
 \norm{\nabla f}_\infty
 \le C_\ell\left(
 R^{3/2}\norm{\Lambda f}_2
 +R^{5/2-\ell}\norm f_{H^\ell}
 \right).
\]
If \(\norm{\Lambda f}_2=0\), then \(f=0\) in \(H^\ell(\R^3)\), and the
estimate is immediate.  Otherwise, balance the two terms by choosing
\(R^{\ell-1}=\norm f_{H^\ell}/\norm{\Lambda f}_2\).  This proves
\eqref{eq:finite-order-fourier-split}.

Let \(T_{\max}\) be the maximal classical lifespan from \(u_0\), and fix
\[
 0<\tau<\min\{T,T_{\max}\}.
\]
Apply the Fourier inequality on \([0,\tau]\), then insert
\eqref{eq:finite-order-lipschitz-persistence}.  Since
\(\lambda_\ell=C_\ell^{\rm Lip}\theta_\ell\), throughout this interval
\[
 -\frac d{dt}e^{-\lambda_\ell\mathcal L(t)}
 \le \lambda_\ell G_\ell
 \norm{u_0}_{H^\ell}^{\theta_\ell}
 \norm{\Lambda u(t)}_2^{1-\theta_\ell}.
\]
The kinetic-energy identity and H\"older's inequality give
\[
 \int_0^\tau\norm{\Lambda u(t)}_2^{1-\theta_\ell}\,dt
 \le \tau^{(1+\theta_\ell)/2}
 \left(\frac{\norm{u_0}_2^2}{2\nu}\right)^{
 (1-\theta_\ell)/2}
 \le T^{(1+\theta_\ell)/2}
 \left(\frac{\norm{u_0}_2^2}{2\nu}\right)^{
 (1-\theta_\ell)/2}.
\]
Integration now yields
\[
 e^{-\lambda_\ell\mathcal L(\tau)}
 \ge1-\Xi_\ell[u_0,\nu;T],
\]
and hence the right side of
\eqref{eq:finite-order-lipschitz-clock} bounds
\(\mathcal L(\tau)\), uniformly in every such \(\tau\).
If \(T_{\max}\le T\), this bound and
\eqref{eq:finite-order-lipschitz-persistence} keep the \(H^\ell\) norm
bounded as \(\tau\uparrow T_{\max}\).  The continuation estimate
\eqref{eq:high-order-gronwall} would then extend the classical solution
past \(T_{\max}\), a contradiction.  Thus \(T<T_{\max}\); letting
\(\tau\uparrow T\) proves \eqref{eq:finite-order-lipschitz-clock}.

It remains to retain the signed critical value rather than only the
Lipschitz clock.  Equation~\eqref{eq:lipschitz-enstrophy-row} gives
\begin{equation*}
 \sup_{0\le t\le T}Y(t)
 \le D_1e^{2C_1^{\rm Lip}\mathcal L(T)},
 \qquad
 2\nu\int_0^TZ(t)\,dt
 \le D_1e^{2C_1^{\rm Lip}\mathcal L(T)}.
\end{equation*}
Energy also gives
\(\int_0^TY(t)\,dt\le E_0/(2\nu)\).  Spatial interpolation on each time
slice, followed by Cauchy--Schwarz in time, now gives
\[
 \nu\int_0^T\norm{\Lambda^{3/2}u(t)}_2^2\,dt
 \le\frac12\sqrt{E_0D_1}\,
 e^{C_1^{\rm Lip}\mathcal L(T)}.
\]
At every \(S\le T\),
\[
 H(S)
 \le\frac12\norm{u(S)}_2\norm{\Lambda u(S)}_2
 \le\frac12\sqrt{E_0D_1}\,
 e^{C_1^{\rm Lip}\mathcal L(T)}.
\]
Insert these two bounds into the exact critical balance and take the supremum
over \(S\).  The balance subtracts \(H(0)\), producing the last term in
\eqref{eq:finite-order-critical-prefix}.  Notice that
\[
 2H(0)=\norm{\Lambda^{1/2}u_0}_2^2
 \le\sqrt{E_0D_1},
\]
so the displayed majorant is nonnegative.  Equation
\eqref{eq:finite-order-lipschitz-clock} converts the last exponential into
the power in
\eqref{eq:finite-order-critical-prefix}, and
\eqref{eq:finite-order-critical-coefficient} follows from its definition.
\end{proof}

Fix an integer \(m\ge1\).  Differentiating the equation \(m\) times in
space, pairing with the same derivative of \(u\), and summing with the
multinomial weights keeps the cancellation of the undifferentiated transport
term and of pressure.  If \(k\) derivatives land on the transporting
velocity, Gagliardo--Nirenberg interpolation places that factor between
\(\nabla u\in L^3\) and \(\nabla^m u\in L^6\); the remaining differentiated
factor uses the complementary interpolation exponent.  H\"older's inequality
then gives the complete-row estimate
\begin{equation}
 |\mathcal P_m(t)|
 \le K_m
 \norm{\Lambda^{3/2}u(t)}_2
 \norm{\Lambda^{m+1}u(t)}_2
 \norm{\Lambda^m u(t)}_2,
 \qquad m\ge1,
 \label{eq:critical-coefficient-row-estimate}
\end{equation}
where \(K_m\) depends only on the derivative order and the fixed Sobolev
constants.  This is the spatially differentiated VPI row: every product term
is retained before the common estimate is taken.

Set
\[
 X_m(t):=\norm{\Lambda^m u(t)}_2^2,
 \qquad
 Y_m(t):=\norm{\Lambda^{m+1}u(t)}_2^2,
 \qquad
 D_m(u_0):=\norm{\Lambda^m u_0}_2^2.
\]
Young's inequality in \eqref{eq:critical-coefficient-row-estimate} gives
\begin{equation*}
 X_m'(t)+\nu Y_m(t)
 \le\frac{K_m^2}{\nu}
 \norm{\Lambda^{3/2}u(t)}_2^2X_m(t).
\end{equation*}
The producer coefficient bound \eqref{eq:producer-critical-coefficient-bound}
controls the complete coefficient.
With
\begin{equation}
 \Gamma_m[u_0,\nu;T_*]
 :=\frac{K_m^2}{\nu}
 \mathfrak A_{1/2}^{\rm prod}[u_0,\nu;T_*]
 \label{eq:spatial-envelope-exponent}
\end{equation}
Gronwall and one subsequent integration give the explicit spatial bounds
\begin{equation}
 \begin{aligned}
  \sup_{0\le t<T_*}\norm{\Lambda^m u(t)}_2^2
  &\le D_m(u_0)e^{\Gamma_m},\\
  \nu\int_0^{T_*}\norm{\Lambda^{m+1}u(t)}_2^2\,dt
  &\le D_m(u_0)e^{\Gamma_m}.
 \end{aligned}
 \label{eq:explicit-spatial-envelope}
\end{equation}
The dependence on the derivative order is carried by \(D_m\) and \(K_m\);
the viscosity appears both in the exponent and in the dissipative factor.

It will be useful to place the inhomogeneous norms into one notation.  Define
\begin{equation}
 \mathfrak S_0:=\norm{u_0}_2,
 \qquad
 \mathfrak S_m^2
 :=2^{m-1}\left(
  \norm{u_0}_2^2+D_m(u_0)e^{\Gamma_m}
 \right)
 \quad(m\ge1).
 \label{eq:spatial-sup-envelope}
\end{equation}
Since \((1+r)^m\le2^{m-1}(1+r^m)\) for \(r\ge0\),
\begin{equation}
 \sup_{0\le t<T_*}\norm{u(t)}_{H^m}\le\mathfrak S_m
 \qquad(m\in\Nzero).
 \label{eq:spatial-sup-envelope-bound}
\end{equation}
The integrated companion begins with the exact energy row,
\begin{equation}
 \mathfrak D_0^2
 :=T_*\norm{u_0}_2^2+\frac{\norm{u_0}_2^2}{2\nu},
 \qquad
 \norm u_{L^2(I_*;H^1)}\le\mathfrak D_0.
 \label{eq:spatial-dissipation-envelope-zero}
\end{equation}
For \(m\ge1\), set
\begin{equation}
 \mathfrak D_m^2
 :=2^m\left[
  T_*\norm{u_0}_2^2
  +\frac{D_m(u_0)}{\nu}e^{\Gamma_m}
 \right],
 \qquad
 \norm u_{L^2(I_*;H^{m+1})}\le\mathfrak D_m.
 \label{eq:spatial-dissipation-envelope}
\end{equation}

The temporal rows now inherit an entirely finite recursion.  Put
\(\sigma_r:=\max\{2,r+1\}\).  Enlarge the order-dependent constants from
Appendix~\ref{app:analytic-framework}, if needed, so that for every integer
\(r\ge0\),
\begin{equation}
 \begin{aligned}
 &\norm{\PP\nabla\cdot(f\otimes g)}_{H^r}
 +\norm{R_iR_j(f_i g_j)}_{H^r}\\
 &\hspace{5em}\le
 A_r\norm f_{H^{\sigma_r}}\norm g_{H^{\sigma_r}}.
 \end{aligned}
 \label{eq:row-product-constant}
\end{equation}
At \(r=1\), this product estimate and
\eqref{eq:spatial-sup-envelope-bound} close the nonlinear graph action:
\begin{equation}
 \mathfrak J_1[u_0,\nu;T_*]
 \le A_1^2\int_0^{T_*}\norm{u(t)}_{H^2}^4\,dt
 \le A_1^2T_*\mathfrak S_2^4.
 \label{eq:critical-current-to-graph-action}
\end{equation}
There is a sharper estimate that uses the differentiated spatial envelopes
instead of the uniform \(H^2\) norm.  Write
\[
 a(t):=\norm{\nabla u(t)}_2,
 \qquad
 b(t):=\norm{\nabla^2u(t)}_2.
\]
Sobolev interpolation and boundedness of \(\PP\) give a constant \(G_1\),
depending only on the fixed three-dimensional inequalities, such that
\begin{equation}
 \norm{B(u,u)}_{H^1}^2
 \le G_1\bigl(a^3b+ab^3\bigr).
 \label{eq:graph-action-product-sharp}
\end{equation}
Indeed,
\[
 \norm{u\cdot\nabla u}_2
 \le C a^{3/2}b^{1/2},
 \qquad
 \norm{\nabla(u\cdot\nabla u)}_2
 \le C a^{1/2}b^{3/2}.
\]
The two time integrals in \eqref{eq:graph-action-product-sharp} retain one
supremum and one copy of the energy--dissipation pairing:
\begin{align*}
 \int_0^{T_*}a^3b\,dt
 &\le \sup_{t<T_*}a(t)^2
       \left(\int_0^{T_*}a^2\,dt\right)^{1/2}
       \left(\int_0^{T_*}b^2\,dt\right)^{1/2},\\
 \int_0^{T_*}ab^3\,dt
 &\le \sup_{t<T_*}b(t)^2
       \left(\int_0^{T_*}a^2\,dt\right)^{1/2}
       \left(\int_0^{T_*}b^2\,dt\right)^{1/2}.
\end{align*}
The kinetic-energy identity and
\eqref{eq:explicit-spatial-envelope} at \(m=1,2\) now yield
\begin{equation}
 \mathfrak J_1[u_0,\nu;T_*]
 \le \mathfrak J_1^{\sharp}[u_0,\nu;T_*]
 :=\frac{G_1}{\sqrt2\,\nu}
 \sqrt{\norm{u_0}_2^2\,D_1(u_0)}\,e^{\Gamma_1/2}
 \left(D_1(u_0)e^{\Gamma_1}+D_2(u_0)e^{\Gamma_2}\right).
 \label{eq:graph-action-critical-sharp}
\end{equation}
This version has no explicit factor of \(T_*\).  Its horizon dependence is
carried by the effective coefficient inside \(\Gamma_1\) and \(\Gamma_2\).
Together, \eqref{eq:critical-current-to-graph-action} and
\eqref{eq:graph-action-critical-sharp} give
\begin{equation}
 \mathfrak J_1[u_0,\nu;T_*]
 \le \mathfrak J_1^{\mathrm{fin}}
 :=\min\left\{
 A_1^2T_*\mathfrak S_2^4,
 \mathfrak J_1^{\sharp}[u_0,\nu;T_*]
 \right\}.
 \label{eq:graph-action-best-finite-bound}
\end{equation}
Substitution into \eqref{eq:base-producer-graph-bound} gives the analytic
feedback estimate
\begin{equation}
 \begin{aligned}
 \mathfrak B_2[u_0,\nu;T_*]
 &\le \overline{\mathfrak B}_2[u_0,\nu;T_*],\\
 \overline{\mathfrak B}_2[u_0,\nu;T_*]
 &:=4T_*\norm{u_0}_2^2
 +\kappa_\nu\left(
   \nu Z_0+\mathfrak J_1^{\mathrm{fin}}
  \right).
 \end{aligned}
 \label{eq:closed-base-rectangle-current-form}
\end{equation}
Here \(\mathfrak S_2\) is the explicit expression
\eqref{eq:spatial-sup-envelope}, with exponent
\eqref{eq:spatial-envelope-exponent}.  The right side of
\eqref{eq:closed-base-rectangle-current-form} depends on
\(u_0,\nu,T_*\), the fixed analytic constants, and the single producer
coefficient \(\mathfrak A_{1/2}^{\rm prod}\).  By
\eqref{eq:base-rectangle-critical-coefficient}, the current-free choice
\(\mathfrak A_{1/2}^{\rm prod}\leq\mathfrak B_2\) reduces every higher
readout to the one base output of the datum producer.  Thus
\(\mathfrak B_2\) is the sole solution-dependent input; no higher rectangle
value remains on the right.  This form exposes the nonlinear feedback but does
not yet remove the base output itself.  After global continuation has been
proved, Section~\ref{sec:final-estimates} removes that final dependence by a
residual certificate computed from the datum.

Starting from \(\mathfrak V_{0,r}:=\mathfrak S_r\), define
\begin{equation}
 \mathfrak V_{n+1,r}
 :=\nu\mathfrak V_{n,r+2}
 +A_r\sum_{a=0}^{n}\binom na
 \mathfrak V_{a,\sigma_r}\mathfrak V_{n-a,\sigma_r},
 \qquad n,r\in\Nzero,
 \label{eq:explicit-velocity-row-recursion}
\end{equation}
and
\begin{equation}
 \mathfrak Q_{n,r}
 :=A_r\sum_{a=0}^{n}\binom na
 \mathfrak V_{a,\sigma_r}\mathfrak V_{n-a,\sigma_r}.
 \label{eq:explicit-pressure-row-recursion}
\end{equation}
The projected momentum recurrence and the normalized Riesz formula give, by
induction in \(n\),
\begin{equation}
 \sup_{0\le t<T_*}\norm{V_n(t)}_{H^r}
 \le\mathfrak V_{n,r},
 \qquad
 \sup_{0\le t<T_*}\norm{Q_n(t)}_{H^r}
 \le\mathfrak Q_{n,r}.
 \label{eq:explicit-jet-row-bounds}
\end{equation}
Every entry on the right of \eqref{eq:explicit-velocity-row-recursion} has
lower temporal order, while the two additional spatial derivatives pay the
Laplacian and the differentiated quadratic term.  The binomial sum retains
all temporal distributions of the original transport product.

The lowest rectangle can keep its natural negative row.  Define
\begin{equation}
 \mathfrak V_{n+1,-1}
 :=\nu\mathfrak V_{n,1}
 +A_{-1}\sum_{a=0}^{n}\binom na
 \mathfrak V_{a,1}\mathfrak V_{n-a,1},
 \label{eq:negative-velocity-row-envelope}
\end{equation}
where \(A_{-1}\) is the constant in
\[
 \norm{\PP\nabla\cdot(f\otimes g)}_{H^{-1}}
 \le A_{-1}\norm f_{H^1}\norm g_{H^1}.
\]
The projected recurrence gives
\(\sup_{t<T_*}\norm{V_{n+1}(t)}_{H^{-1}}
\le\mathfrak V_{n+1,-1}\).

For \(M\in\Nzero\), set
\begin{equation}
 \mathfrak R_{N,M}[u_0,\nu;T_*]
 :=T_*\sum_{n=0}^{N}\left(
  \mathfrak V_{n,M+1}^2
  +\mathfrak V_{n+1,M-1}^2
 \right).
 \label{eq:explicit-rectangle-envelope}
\end{equation}
For \(M=0\), \eqref{eq:negative-velocity-row-envelope} controls the lower
row.  For \(M\ge1\), the nonnegative recursion controls it.  In either case,
\begin{equation}
 \mathcal N_{N,M}(u;I_*)
 \le\mathfrak R_{N,M}[u_0,\nu;T_*]<\infty.
 \label{eq:closed-rectangle-estimate}
\end{equation}
Once \(N\) and \(M\) are chosen, the right side is a finite recursive
expression in \(\mathfrak A_{1/2}^{\rm prod}[u_0,\nu;T_*]\), \(\nu\), \(T_*\), the named
analytic constants through the required orders, and the corresponding finite
Sobolev norms of \(u_0\).  The velocity rectangle norm \((N,M)\) uses datum
derivatives only through order \(M+2N+1\).  If the retained pressure row
\(Q_N\in H^M\) is included, the current
\(\sigma_M=\max\{2,M+1\}\) recursion also uses order
\(\sigma_M+2N\): this is \(M+2N+1\) for \(M\ge1\) and \(2N+2\) for
\(M=0\).  No higher rectangle value remains on the right.

\Needspace{18\baselineskip}
\subsection{Classical continuation quantities as finite projections}

The Ladyzhenskaya--Prodi--Serrin scale is the first direct example.  Select the
zeroth temporal row in \(L^2(I_*;H^2)\).  Since
\(H^2(\R^3)\hookrightarrow L^\infty(\R^3)\),
\[
 u\in L^2(I_*;L^\infty),
 \qquad
 \frac22+\frac3\infty=1.
\]
One finite coordinate has reached the scaling line.  More generally, if
\(b>3\) and \(1\le a\le\infty\) satisfy
\[
 \frac2a+\frac3b\le1,
\]
equation~\eqref{eq:finite-projection-mixed-norms} with \(n=k=0\) gives
 \[
 u\in L^a(I_*;L^b).
 \]
More precisely, for any integer
\(m>\frac32-\frac3b\),
\begin{equation}
 \norm{u}_{L^a(I_*;L^b)}
 \le C_{m,b}T_*^{1/a}\mathfrak S_m.
 \label{eq:quantitative-lps-family}
\end{equation}
Thus every classical Ladyzhenskaya--Prodi--Serrin pair with \(b>3\) is a
finite projection of the same compatible family.

The same calculation gives the final mixed-jet estimates at arbitrary finite
orders.  If \(m>k+\frac32-\frac3b\), then
\begin{equation}
 \begin{aligned}
  \sup_{0\le t<T_*}
  \norm{\nabla^k\partial_t^n u(t)}_{L^b}
  &\le C_{m,k,b}\mathfrak V_{n,m},\\
  \norm{\nabla^k\partial_t^n u}_{L^a(I_*;L^b)}
  &\le C_{m,k,b}T_*^{1/a}\mathfrak V_{n,m}.
 \end{aligned}
 \label{eq:closed-mixed-jet-estimates}
\end{equation}
If \(s\in\Nzero\) satisfies
\(s>k+\frac32-\frac3b\), the normalized pressure rows obey
\begin{equation}
 \begin{aligned}
  \sup_{0\le t<T_*}\norm{\nabla^kQ_n(t)}_{L^b}
  &\le C_{s,k,b}\mathfrak Q_{n,s},\\
  \norm{\nabla^kQ_n}_{L^a(I_*;L^b)}
  &\le C_{s,k,b}T_*^{1/a}\mathfrak Q_{n,s}.
 \end{aligned}
 \label{eq:closed-pressure-jet-estimates}
\end{equation}
Here and below \(T_*^{1/\infty}=1\).

The same \(L^2(I_*;H^3)\) coordinate controls the stretching quantities from
Part~I.  Put
\[
 \omega:=\nabla\times u,
 \qquad
 S:=\frac12\bigl(\nabla u+\nabla u^{\mathsf T}\bigr).
\]
Spatial embedding gives
\[
 \omega,\ S,\ \nabla u\in L^2(I_*;L^\infty).
\]
Since \(T_*<\infty\), Cauchy--Schwarz in time gives
\begin{equation}
 \int_0^{T_*}
 \bigl(\norm{\omega(t)}_\infty+\norm{S(t)}_\infty
       +\norm{\nabla u(t)}_\infty\bigr)\,dt
 \le C T_*^{1/2}\mathfrak B_2[u_0,\nu;T_*]^{1/2}
 \le C T_*^{1/2}\overline{\mathfrak B}_2[u_0,\nu;T_*]^{1/2}<\infty.
 \label{eq:quantitative-stretching-actions}
\end{equation}
The vorticity and strain actions are not new assumptions at the end of the
proof.  They are low spatial projections of the rectangle that already
produced \eqref{eq:terminal-lipschitz-bound}.

Pressure is another projection of the same object.  The normalized identity
\[
 p=R_iR_j(u_i u_j)
\]
and the pressure projection at the end of Part~III give
\(p\in C([0,T_*];H^m)\) for every fixed \(m\).  Choosing \(m=3\) and using
\eqref{eq:appendix-spatial-embedding} gives
\begin{equation}
 \int_0^{T_*}\norm{\nabla p(t)}_\infty\,dt
 \le C\norm p_{L^1(I_*;H^3)}
 \le C\norm u_{L^2(I_*;H^4)}^2
 \le C\mathfrak R_{0,3}[u_0,\nu;T_*]<\infty.
 \label{eq:quantitative-pressure-action}
\end{equation}
The same argument applies to any fixed temporal pressure row after a finite
spatial margin is added.

The critical endpoint is a second direct projection.
Equation~\eqref{eq:all-row-traces} with \(n=0\) and \(m=1\) gives
\[
 u\in C([0,T_*];H^1).
\]
Since \(H^1\hookrightarrow L^6\) and \(H^1\subset L^2\), interpolation gives
\[
 \norm{u(t)}_3
 \le \norm{u(t)}_2^{1/2}\norm{u(t)}_6^{1/2}
 \le C\norm{u(t)}_{H^1}.
\]
Thus
\begin{equation}
 \sup_{0\le t<T_*}\norm{u(t)}_3
 \le C\mathfrak S_1<\infty,
 \label{eq:quantitative-L3-endpoint}
\end{equation}
which is incompatible with the necessary \(L^3\) limsup behavior in the
Escauriaza--Seregin--\v{S}ver\'ak theorem\bibref{ESS}.  It also excludes the
full terminal divergence
\eqref{eq:seregin-L3-divergence} when \(u_0\) belongs to the compactly supported
subclass used for Seregin's full terminal theorem\bibref{Seregin2012}.

The \(H^1\) trace that bounds the \(L^3\) endpoint also controls the critical
height that generated the response tower in Part~II.  Split frequency at
\(|\xi|=1\):
\begin{align*}
 2H(t)
 &=\int_{|\xi|\le1}|\xi|\,|\widehat u(\xi,t)|^2\,d\xi
 +\int_{|\xi|>1}|\xi|\,|\widehat u(\xi,t)|^2\,d\xi\\
 &\le \norm{u(t)}_2^2+\norm{\nabla u(t)}_2^2
 \le \norm{u(t)}_{H^1}^2.
\end{align*}
Hence
\begin{equation}
 H_{\max}:=\sup_{0\le t<T_*}H(t)
 \le\frac12\mathfrak S_1^2<\infty.
 \label{eq:terminal-critical-height-bound}
\end{equation}
Let \(N_{\mathrm{pass}}\) denote the
number of completed passages in the dyadic ladder from Part~II.  If
\(N_{\mathrm{pass}}\ge1\), the last endpoint level reached is
\[
 L_{N_{\mathrm{pass}}+1}
 =2^{N_{\mathrm{pass}}}L_1,
\]
and hence
\[
 2^{N_{\mathrm{pass}}}L_1\le H_{\max}.
\]
Thus \(N_{\mathrm{pass}}=0\) when \(H_{\max}<L_1\), while
\[
 N_{\mathrm{pass}}
 \le\left\lfloor\log_2\frac{H_{\max}}{L_1}\right\rfloor
 \qquad(H_{\max}\ge L_1).
\]
The height bound does not say that the solution never rises.  It says that the
infinite terminal ladder demanded by \eqref{eq:critical-height-limsup} cannot
be completed.
For the compactly supported subclass used in Part~I, this also excludes the
necessary divergence in Seregin's critical Sobolev endpoint
theorem\bibref{SereginCritical}.  Both Seregin comparisons concern the
narrower compactly supported data class, while the compatible-family argument
uses the stated Schwartz data.

\subsection{The first rectangle seen through the critical height}

The height bound ends the level ladder by bounding \(H\).  The same
\((0,2)\) rectangle also resolves the mechanism inside each attempted
passage: the signed response and the next half-integer dissipation from
Part~II.  To expose that critical block on the complete terminal interval,
retain both adjacent rows:
\begin{equation*}
 u\in L^2(I_*;H^3),
 \qquad
 u_t\in L^2(I_*;H^1).
\end{equation*}
In particular, \(u\in L^2H^{5/2}\) and
\(u_t\in L^2H^{1/2}\).  Lemma~\ref{lem:part-ii-adjacent-trace} at
\(m=3/2\) gives
\[
 u\in C([0,T_*];H^{3/2}).
\]

This pair also controls the variation of the next critical energy.  Let
\(S_R\) be a smooth self-adjoint Fourier cutoff, equal to one on
\(|\xi|\le R\) and supported in \(|\xi|\le2R\), and put \(u_R=S_Ru\).
For classical times,
\begin{align*}
 \frac d{dt}\frac12\norm{\Lambda^{3/2}u_R(t)}_2^2
 &=\operatorname{Re}
   (\Lambda^{3/2}S_Ru_t,\Lambda^{3/2}S_Ru)_2\\
 &=\operatorname{Re}
   (\Lambda^{1/2}S_Ru_t,\Lambda^{5/2}S_Ru)_2.
\end{align*}
Integrate from \(r\) to \(t\).  The two cutoff factors converge in
\(L^2(I_*;L^2)\), so their product converges in \(L^1(I_*)\).  The
\(H^{3/2}\) trace handles the endpoint energies.  Passing \(R\to\infty\)
gives
\begin{equation*}
 E_{3/2}(t)-E_{3/2}(r)
 =\int_r^t\operatorname{Re}
 (\Lambda^{1/2}u_\tau(\tau),\Lambda^{5/2}u(\tau))_2\,d\tau.
\end{equation*}
The integrand is absolutely integrable because
\begin{equation*}
 \int_0^{T_*}
 \left|(\Lambda^{1/2}u_t,\Lambda^{5/2}u)_2\right|dt
 \le
 \norm{u_t}_{L^2(I_*;H^{1/2})}
 \norm{u}_{L^2(I_*;H^{5/2})}
 \le\frac12\mathcal N_{0,2}(u;I_*).
\end{equation*}
Moreover,
\begin{equation}
 \int_0^{T_*}E_{3/2}(t)\,dt
 \le\frac12\norm{u}_{L^2(I_*;H^3)}^2
 \le\frac12\mathfrak B_2[u_0,\nu;T_*]
 \le\frac12\overline{\mathfrak B}_2[u_0,\nu;T_*],
 \label{eq:quantitative-critical-dissipation}
\end{equation}
while \eqref{eq:rectangle-to-trace} at the middle order \(2\) gives
\begin{equation}
 \sup_{0\le t\le T_*}E_{3/2}(t)
 \le\frac12\mathfrak S_2^2.
 \label{eq:quantitative-next-critical-height}
\end{equation}
Thus
\begin{equation*}
 E_{3/2}\in W^{1,1}(I_*)\cap L^1(I_*)
 \subset C([0,T_*])\cap L^\infty(I_*).
\end{equation*}

A second cutoff calculation at order \(1/2\) gives
\[
 H'(t)=\operatorname{Re}
 (\Lambda^{1/2}u_t(t),\Lambda^{1/2}u(t))_2
\]
in distributions.  Both factors are square-integrable in time, and the
zeroth spatial row gives \(H\in L^1(I_*)\).  Thus
\begin{equation}
 H\in W^{1,1}(I_*)\hookrightarrow C([0,T_*]),
 \qquad
 \norm{H'}_{L^1(I_*)}
 \le\frac12\mathfrak B_2[u_0,\nu;T_*]
 \le\frac12\overline{\mathfrak B}_2[u_0,\nu;T_*].
 \label{eq:quantitative-critical-variation}
\end{equation}
The exact critical balance \eqref{eq:H1} now yields
\begin{equation}
 \mathcal P_{1/2}=H'+2\nu E_{3/2}\in L^1(I_*),
 \qquad
 \norm{\mathcal P_{1/2}}_{L^1(I_*)}
 \le\left(\frac12+\nu\right)\mathfrak B_2[u_0,\nu;T_*]
 \le\left(\frac12+\nu\right)
      \overline{\mathfrak B}_2[u_0,\nu;T_*].
 \label{eq:quantitative-critical-response}
\end{equation}
The complete first critical block is
\begin{equation}
 H\in W^{1,1}(I_*),
 \qquad
 E_{3/2}\in L^\infty(I_*)\cap L^1(I_*),
 \qquad
 \mathcal P_{1/2}\in L^1(I_*).
 \label{eq:critical-block}
\end{equation}
Thus \eqref{eq:critical-block} records finite critical variation, integrable
next dissipation, and the complete signed response on the terminal interval.

The same spatial-row estimate controls every fixed height readout.  For
\(j\in\Nzero\), the complete row gives
\(u\in L^2(I_*;H^{j+1})\).  Since
\(|\xi|^{2j+1}\le(1+|\xi|^2)^{j+1}\),
\begin{align*}
 \int_0^{T_*}E_{j+1/2}(t)\,dt
 &=\frac12\int_0^{T_*}
 \norm{\Lambda^{j+1/2}u(t)}_2^2\,dt
 \\
 &\le\frac12\norm{u}_{L^2(I_*;H^{j+1})}^2
 \le\frac12\mathfrak D_j^2<\infty.
\end{align*}
By Lemma~\ref{lem:height-tower}, this is the spatial energy exposed
by the \(j\)-th differentiated critical-height row.  Temporal ascent and
spatial depth remain different coordinates, but every fixed spatial readout
they expose is controlled on the same terminal interval.

\subsection{Uniform tightness of the critical tail}

The finite height bound stops the level ladder.  Part~II's outward scale
migration leaves a second question: can the activity responsible for the
terminal approach retreat through an unbounded collection of higher
frequencies?  Part~I identified vorticity as the local stretching diagnostic.
When that activity is distributed across the shell index, its complete dyadic
\(L^\infty\) sum is the relevant quantity.  The first estimate controls this
sum in time; the second returns to \(H\) and makes its critical tail vanish
uniformly.

For the inhomogeneous decomposition from Appendix~A, Bernstein's inequality
gives, for \(j\ge0\),
\[
 \norm{P_j\omega(t)}_\infty
 \le C2^{3j/2}\norm{P_j\omega(t)}_2
 =C2^{-j/2}\bigl(2^{2j}\norm{P_j\omega(t)}_2\bigr).
\]
Cauchy--Schwarz in \(j\), followed by the dyadic Sobolev equivalence, gives
\[
 \sum_{j\ge0}\norm{P_j\omega(t)}_\infty
 \le C\norm{\omega(t)}_{H^2}
 \le C\norm{u(t)}_{H^3}.
\]
The low block \(S_{1/2}\omega\) satisfies the same estimate by the
low-frequency Bernstein bound.  Integration over the finite interval yields
\begin{align*}
 \int_0^{T_*}
 \left(
  \norm{S_{1/2}\omega(t)}_\infty
  +\sum_{j\ge0}\norm{P_j\omega(t)}_\infty
 \right)\,dt
 &\le CT_*^{1/2}\norm{u}_{L^2(I_*;H^3)} \\
 &\le CT_*^{1/2}\mathfrak B_2[u_0,\nu;T_*]^{1/2} \\
 &\le CT_*^{1/2}\overline{\mathfrak B}_2[u_0,\nu;T_*]^{1/2}<\infty.
\end{align*}
This is how the intersection controls a summation over scales: a fixed shell
weight is dominated by one sufficiently high finite Sobolev coordinate.
More generally, fix \(\sigma\in\R\) and \(k,n\in\Nzero\), and choose an
integer \(M\ge1\) satisfying
\[
 M>\sigma+k+\frac32.
\]
Bernstein and Cauchy--Schwarz in the shell index give
\[
 \norm{S_{1/2}\nabla^k\partial_t^n u(t)}_\infty
 +\sum_{j\ge0}2^{\sigma j}
  \norm{P_j\nabla^k\partial_t^n u(t)}_\infty
 \le C_{\sigma,k,M}\norm{\partial_t^n u(t)}_{H^M}.
\]
Selecting the corresponding finite rectangle and integrating in time yields
\begin{align*}
 \int_0^{T_*}
 \left(
  \norm{S_{1/2}\nabla^k\partial_t^n u(t)}_\infty
  +\sum_{j\ge0}2^{\sigma j}
   \norm{P_j\nabla^k\partial_t^n u(t)}_\infty
 \right)dt
 &\le C_{\sigma,k,M}T_*
 \mathfrak V_{n,M}<\infty.
\end{align*}
Thus every fixed dyadic weight, spatial derivative, and temporal row is
supplied by one sufficiently high finite coordinate.

Fix an integer \(m\ge1\).  For \(K\ge1\), define
\[
 H_{>K}(t)
 :=\frac12\int_{|\xi|>K}|\xi|\,|\widehat u(t,\xi)|^2\,d\xi.
\]
On \(|\xi|>K\),
\(|\xi|\le K^{1-2m}|\xi|^{2m}\), so
\begin{equation}
 \sup_{0\le t<T_*}H_{>K}(t)
 \le\frac12K^{1-2m}
 \sup_{0\le t<T_*}\norm{u(t)}_{\dot H^m}^2
 \le\frac12K^{1-2m}D_m(u_0)e^{\Gamma_m}
 \longrightarrow0
 \quad(K\to\infty).
 \label{eq:uniform-critical-tail}
\end{equation}
The convergence in \eqref{eq:uniform-critical-tail} is uniform in time along
the proposed terminal approach.
Consequently, any critical height carried beyond frequency \(K\) vanishes
uniformly as \(K\to\infty\).  This is the scale-level conclusion behind the
failure of the Zeno scenario: the whole-terminal estimate removes both its
unbounded height sequence and its ultraviolet residue.

Thus the critical-height scenario from Part~II cannot occur.  It required a
rising scale-invariant norm, shorter level passages, and energy pushed toward
higher frequencies.  Whole-terminal control gives the opposite conclusion:
every selected finite coordinate is bounded, every fixed critical tail
vanishes uniformly, and the sequence of higher level passages cannot continue
through infinitely many levels.

This excludes Blown at every fixed finite coordinate, but it has not yet made
the separately obtained traces into one state.  A list of finite limits is not
enough: compatibility must identify one velocity jet across all Sobolev
topologies, and the limiting rows must still satisfy the normalized pressure
and momentum laws.  The next two sections pay those remaining obligations.

\section{Every temporal row reaches one endpoint}
\label{sec:terminal-row-traces}

The projection argument in \eqref{eq:all-row-traces} has carried each
temporal row to the endpoint in every fixed Sobolev space.  The readouts above
used those bounds one target at a time.  The new point is identity: the
\(H^m\) value and the \(H^k\) value obtained from different finite rectangles
must belong to one endpoint jet of the original history.

Write \(V_n^{*,m}=V_n(T_*)\in H^m\).  If \(k>m\), the
\(H^k\)-continuous and \(H^m\)-continuous representatives both equal the
classical row on \(I_*\).  After the embedding \(H^k\hookrightarrow H^m\),
uniqueness of the continuous representative gives
\[
 V_n^{*,k}=V_n^{*,m}\quad\hbox{in }H^m.
\]
We can now drop the topology from the notation and set
\begin{equation}
 V_n^*:=\lim_{t\uparrow T_*}V_n(t)
 \in\bigcap_{m=0}^{\infty}H^m(\R^3)=H^\infty(\R^3).
 \label{eq:all-temporal-endpoints}
\end{equation}
The trace retains the size of the rectangle that produced it.  For every
\(n,m\in\Nzero\),
\begin{equation}
 \norm{V_n^*}_{H^m}
 \le\mathfrak V_{n,m},
 \qquad
 \norm{V_n(t)-V_n^*}_{H^m}
 \le(T_*-t)\mathfrak V_{n+1,m}.
 \label{eq:quantitative-endpoint-row}
\end{equation}
The first estimate follows by passing the uniform row bound
\eqref{eq:explicit-jet-row-bounds} to the trace.  For the second, the
Banach-valued fundamental theorem of calculus gives
\[
 V_n^*-V_n(t)=\int_t^{T_*}V_{n+1}(\tau)\,d\tau
 \quad\hbox{in }H^m,
\]
and the uniform bound for the next temporal row gives the stated Lipschitz
modulus.
For \(n=0\), define
\begin{equation}
 u_*:=V_0^*=\lim_{t\uparrow T_*}u(t)
 \in\bigcap_{m=0}^{\infty}H^m(\R^3)=H^\infty(\R^3).
 \label{eq:common-velocity-endpoint}
\end{equation}
Continuity of divergence from \(H^{m+1}\) to \(H^m\) passes
\(\nabla\cdot u=0\) to this limit.  The \(H^0\) trace and the energy identity
give finite kinetic energy.  In particular,
\[
 u_*\in H^3_\sigma(\R^3),
\]
so the endpoint is admissible initial data for the local construction in
Proposition~\ref{prop:maximal-history}.

The zeroth trace is enough to start a right-hand solution.  The higher traces
answer a different question: whether that solution meets the old one without
a corner in time.  The equation will settle this by showing that the left and
right jets are generated from the same \(u_*\).

Every mixed velocity coordinate reaches the endpoint as well.  For a spatial
multi-index \(\alpha\), choose \(m\ge|\alpha|\).  Continuity of
\(\partial_x^\alpha:H^m\to H^{m-|\alpha|}\) gives
\begin{equation*}
 \partial_x^\alpha V_n(t)
 \longrightarrow \partial_x^\alpha V_n^*
 \quad\hbox{in every fixed admissible Sobolev topology}.
\end{equation*}
Sobolev embedding converts sufficiently high source orders into convergence
in every fixed \(C_b^j\) topology.  The remaining endpoint work belongs to the
normalized pressure and the differentiated momentum equation.

\section{The pressure-complete endpoint jet}
\label{sec:endpoint-jet}

The traces in \eqref{eq:all-temporal-endpoints} carry the nonlinear terms to
\(T_*\), one fixed target topology at a time.  Fix \(q\in\Nzero\) and choose
an integer \(\ell\ge q+3\).  Sobolev
multiplication gives
\[
 \norm{fg}_{H^{q+1}}
 \le C_{q,\ell}\norm f_{H^\ell}\norm g_{H^\ell}.
\]
If \(f(t)\to f_*\) and \(g(t)\to g_*\) in \(H^\ell\), then
\begin{align*}
 \norm{f(t)g(t)-f_*g_*}_{H^{q+1}}
 &\le C_{q,\ell}
 \norm{f(t)-f_*}_{H^\ell}\norm{g(t)}_{H^\ell}\\
 &\quad+C_{q,\ell}
 \norm{f_*}_{H^\ell}\norm{g(t)-g_*}_{H^\ell}
 \longrightarrow0.
\end{align*}
The convergent factor \(g(t)\) is bounded near \(T_*\).  After one spatial
derivative,
\[
 (f(t)\cdot\nabla)g(t)
 \longrightarrow(f_*\cdot\nabla)g_*
 \quad\hbox{in }H^q.
\]
Riesz transforms preserve this convergence in every finite Sobolev space.

For each \(n\in\Nzero\), define the normalized terminal pressure row by
\begin{equation}
 Q_n^*
 :=R_iR_j\sum_{a=0}^{n}\binom na
 V_{a,i}^*V_{n-a,j}^*.
 \label{eq:endpoint-pressure}
\end{equation}
Every factor belongs to \(H^\infty\), so \(Q_n^*\in H^\infty\).
If \(s\in\Nzero\) and \(r_s=\max\{s,2\}\), the same product calculation as
\eqref{eq:quantitative-pressure-trace} gives the endpoint bound
\begin{equation}
 \norm{Q_n^*}_{H^s}
 \le\mathfrak Q_{n,s}.
 \label{eq:quantitative-endpoint-pressure}
\end{equation}
The pressure approaches this value with an explicit modulus as well.  With
\(\sigma_s=\max\{2,s+1\}\), put
\begin{equation}
 \begin{aligned}
 \mathfrak L^Q_{n,s}
 :=A_s\sum_{a=0}^{n}\binom na
 \bigl(&
 \mathfrak V_{a+1,\sigma_s}\mathfrak V_{n-a,\sigma_s}\\
 &+\mathfrak V_{a,\sigma_s}
   \mathfrak V_{n-a+1,\sigma_s}
 \bigr).
 \end{aligned}
 \label{eq:pressure-endpoint-modulus-constant}
\end{equation}
Adding and subtracting
\(V_a^*V_{n-a}(t)\) in each product and using
\eqref{eq:quantitative-endpoint-row} gives
\begin{equation}
 \norm{Q_n(t)-Q_n^*}_{H^s}
 \le(T_*-t)\mathfrak L^Q_{n,s}.
 \label{eq:quantitative-pressure-endpoint-modulus}
\end{equation}

\Needspace{14\baselineskip}
\begin{proposition}[Endpoint jet reconstruction]
\label{prop:endpoint-jet}
For every \(n\in\Nzero\), the left velocity traces satisfy
\begin{equation}
 V_{n+1}^*
 =\nu\Delta V_n^*
 -\sum_{a=0}^{n}\binom na
 (V_a^*\cdot\nabla)V_{n-a}^*
 -\nabla Q_n^*.
 \label{eq:endpoint-velocity}
\end{equation}
For every finite \(q\),
\begin{equation}
 \lim_{t\uparrow T_*}
 \left(
  \norm{Q_n(t)-Q_n^*}_{H^q}
  +\norm{V_{n+1}(t)-V_{n+1}^*}_{H^q}
 \right)=0.
 \label{eq:endpoint-jet-convergence}
\end{equation}
The family \((V_n^*,Q_n^*)_{n\ge0}\) is the unique normalized terminal jet
rooted at \(u_*\).
\end{proposition}

\begin{proof}
Fix \(n,q\) and choose \(\ell\ge q+3\).  At every \(t<T_*\),
\[
 Q_n(t)=R_iR_j\sum_{a=0}^{n}\binom na
 V_{a,i}(t)V_{n-a,j}(t).
\]
Equation \eqref{eq:all-row-traces} gives
\(V_a(t)\to V_a^*\) in \(H^\ell\) for every \(a\le n\).  The product
estimate above, applied term by term and followed by the Riesz transforms,
gives
\[
 Q_n(t)\longrightarrow Q_n^*
 \quad\hbox{in }H^{q+1}.
\]
In particular, \(\nabla Q_n(t)\to\nabla Q_n^*\) in \(H^q\).

The velocity recurrence on \(I_*\) is
\[
 V_{n+1}(t)
 =\nu\Delta V_n(t)
 -\sum_{a=0}^{n}\binom na
 (V_a(t)\cdot\nabla)V_{n-a}(t)
 -\nabla Q_n(t).
\]
The trace of \(V_n\) in \(H^{q+2}\) handles diffusion.  The traces in
\(H^\ell\) handle every transport product, and the pressure convergence
handles the final term.  The right side converges in \(H^q\) to the right
side of \eqref{eq:endpoint-velocity}.  At the same time, the adjacent trace
for row \(n+1\) gives
\(V_{n+1}(t)\to V_{n+1}^*\) in \(H^q\).  Uniqueness of limits proves
\eqref{eq:endpoint-velocity} and \eqref{eq:endpoint-jet-convergence}.

Finally, \(V_0^*=u_*\) fixes \(Q_0^*\) through
\eqref{eq:endpoint-pressure}, and \eqref{eq:endpoint-velocity} then fixes
\(V_1^*\).  Those rows fix \(Q_1^*\), which fixes \(V_2^*\), and induction
continues through every finite order.  This proves uniqueness of the normalized
jet.  Spatial differentiation of the convergent identities recovers every
mixed endpoint coordinate.
\end{proof}

The three endpoint failures from Part~I have now been separated and paid.
The finite projections exclude Blown.  Equality of the traces across Sobolev
embeddings excludes Jump.  The limiting pressure and momentum recurrences
exclude Dead.  What remains is not another estimate.  The recovered state has
to start the same equation on the right and meet the old history without a
corner in time.

The terminal rows still satisfy the same incompressible, pressure-normalized
momentum recurrence as the preterminal solution.  Nothing new is chosen for
pressure or for the higher time derivatives.  A right-hand solution starting
from \(u_*\) generates its jet by this same recurrence, so the two jets can now
be compared row by row.

\Needspace{22\baselineskip}
\section{The exact local restart}
\label{sec:restart}

Because \(u_*\in H^\infty_\sigma\), the local theorem applies at \(T_*\) and
produces a unique pressure-normalized solution for a time
\begin{equation}
 \begin{aligned}
 \delta
 &\ge\min\left\{1,c_\nu(1+\norm{u_*}_{H^3})^{-2}\right\}\\
 &\ge\min\left\{
 1,\,
 c_\nu\left(1+\mathfrak S_3\right)^{-2}
 \right\}>0.
 \end{aligned}
 \label{eq:quantitative-terminal-restart}
\end{equation}
Persistence retains every higher spatial order on that same interval, while
the pressure normalization and differentiated equation determine the complete
right-hand jet from \(u_*\).  Proposition~\ref{prop:maximal-history} states
this construction, and Appendix~\ref{app:local-theory} derives its lifespan,
pressure, persistence, and uniqueness in the strong class.

\begin{theorem}[Continuation from the compatible intersection]
\label{thm:intersection-continuation}
Let \((u,p)\) be a pressure-normalized maximal classical solution on
\([0,T_*)\), where \(T_*<\infty\).  If its datum-generated rows satisfy
\eqref{eq:datum-adjacent-row} for every finite \(n,m\), then there are
\(\delta>0\) and a unique pressure-normalized classical solution
\((\widetilde u,\widetilde p)\) on \([0,T_*+\delta]\) whose restriction to
\([0,T_*)\) is \((u,p)\).  The continuation time can be chosen to satisfy
\eqref{eq:quantitative-terminal-restart}.
\end{theorem}

\begin{proof}
Apply Proposition~\ref{prop:maximal-history} at \(t_0=T_*\) with the datum
\(u_*\) constructed in \eqref{eq:common-velocity-endpoint}.  It gives a
unique normalized solution \((w,q)\) on
\([T_*,T_*+\delta]\), where
 \[
  \delta=\delta(\nu,\norm{u_*}_{H^3})>0,
  \qquad w(T_*)=u_*.
 \]
Estimate \eqref{eq:local-lifespan}, followed by
\eqref{eq:spatial-sup-envelope-bound} at \(m=3\), gives
\eqref{eq:quantitative-terminal-restart}.
Because \(u_*\in H^\infty_\sigma\), persistence keeps every finite spatial
Sobolev order on this one interval.  The pressure formula and momentum
recurrence then give
\[
 w,q\in\bigcap_{r,m\in\Nzero}
 C^r([T_*,T_*+\delta];H^m).
\]

Write the right-hand jet as
\[
 W_n^+:=\partial_t^nw(T_*),
 \qquad
 P_n^+:=\partial_t^nq(T_*).
\]
The local solution uses the same pressure normalization, so its initial jet
obeys
\begin{align*}
 P_n^+
 &=R_iR_j\sum_{a=0}^{n}\binom naW_{a,i}^+W_{n-a,j}^+,\\
 W_{n+1}^+
 &=\nu\Delta W_n^+
 -\sum_{a=0}^{n}\binom na(W_a^+\cdot\nabla)W_{n-a}^+
 -\nabla P_n^+.
\end{align*}
The equality \(W_0^+=w(T_*)=u_*=V_0^*\) starts an induction.  If
\(W_a^+=V_a^*\) for \(0\le a\le n\), then the two pressure formulas give
\(P_n^+=Q_n^*\).  The two velocity recurrences then give
\(W_{n+1}^+=V_{n+1}^*\).  Thus
\begin{equation}
 W_n^+=V_n^*,
 \qquad
 P_n^+=Q_n^*
 \qquad(n\in\Nzero).
 \label{eq:left-right-jet-match}
\end{equation}
Continuous spatial differentiation matches every mixed derivative at the
joining time.

The left side already has the boundary regularity needed for gluing.  Part~III
gave
\[
 u\in\bigcap_{r,m\in\Nzero}H^r(I_*;H^m).
\]
For each finite \(a,m\), apply the time-Sobolev embedding
\eqref{eq:appendix-time-embedding} with source order \(a+1\).  This gives
\[
 u\in\bigcap_{a,m\in\Nzero}C^a([0,T_*];H^m),
 \qquad \partial_t^nu(T_*)=V_n^*.
\]
The pressure formula, the temporal Leibniz rule, and product continuity give
the analogous left extension for \(p\), with
\(\partial_t^np(T_*)=Q_n^*\).  Thus both halves are smooth curves into every
fixed \(H^m\), and \eqref{eq:left-right-jet-match} identifies all of their
one-sided derivatives.

Define
\[
 (\widetilde u,\widetilde p)(t)
 :=
 \begin{cases}
  (u,p)(t),&0\le t<T_*,\\
  (w,q)(t),&T_*\le t\le T_*+\delta.
 \end{cases}
\]
For a fixed \(m\), equality at order zero makes the piecewise \(H^m\)-valued
curve continuous.  If it is \(C^k\), equality of the \((k+1)\)-st one-sided
derivatives makes the difference quotient of its \(k\)-th derivative agree
from both sides.  Induction in \(k\) gives
\[
 \widetilde u,\widetilde p
 \in\bigcap_{r,m\in\Nzero}C^r([0,T_*+\delta];H^m).
\]
Taking \(m\) above each spatial embedding threshold gives
\[
 \widetilde u,\widetilde p
 \in C^\infty(\R^3\times[0,T_*+\delta]).
\]

The joined pair solves Navier--Stokes on each side of \(T_*\).  Every term is
continuous there in sufficiently high finite Sobolev spaces, so the equation
also holds at the joining time.  Divergence and the Riesz-transform pressure
normalization pass in the same way.  The joined pair is a classical
normalized solution on a strictly longer interval.  Its left restriction is
the original history by construction, and local uniqueness fixes the
right-hand continuation issued from \(u_*\).
\end{proof}

\Needspace{26\baselineskip}
\section{Global conclusion}

\begin{proof}[Proof of Theorem~\ref{thm:main}]
Fix \(\nu>0\) and \(u_0\in\SSigma(\R^3)\), and let \((u,p)\) be the unique
pressure-normalized maximal classical solution supplied by
Proposition~\ref{prop:maximal-history}, with lifespan \([0,T_*)\).

No finite value of \(T_*\) can be maximal.  At such an endpoint,
Theorem~\ref{thm:whole-terminal} places every finite pressure-complete
rectangle on \((0,T_*)\).  The projection rule gives one terminal value for
every temporal row in every finite Sobolev space, compatibility identifies
those values as one \(H^\infty\) endpoint jet, and
Proposition~\ref{prop:endpoint-jet} shows that the jet still obeys the
normalized momentum and pressure recurrences.
Theorem~\ref{thm:intersection-continuation} then extends the original solution
to a strictly longer interval.  Hence
\[
 T_*=\infty.
\]

Every finite interval now lies inside this one maximal solution.  Persistence
of all spatial orders, followed by the velocity and pressure recurrences, gives
\[
 u,p\in C^\infty(\R^3\times[0,\infty)),
 \qquad
 p(t)=R_iR_j\bigl(u_i(t)u_j(t)\bigr).
\]
The force is zero, \(\nabla\cdot u=0\), and \(u(0)=u_0\).  Finally, the
energy identity \eqref{eq:energy} holds for every \(t\ge0\):
\[
 \frac12\norm{u(t)}_2^2
 +\nu\int_0^t\norm{\nabla u(\tau)}_2^2\,d\tau
 =\frac12\norm{u_0}_2^2.
\]
This is the uniform kinetic-energy bound
\eqref{eq:uniform-kinetic-energy}.
\end{proof}

\Needspace{24\baselineskip}
\section{The final finite-horizon estimates}
\label{sec:final-estimates}

The contradiction argument has ended, so the constants can now be stated on
an arbitrary physical horizon.  Fix \(T>0\), write \(I_T=(0,T)\), and begin
with the one quantity through which all of the preceding finite recursions
run:
\begin{equation}
 \mathfrak B_2[u_0,\nu;T]
 :=\norm u_{L^2(I_T;H^3)}^2
   +\norm{u_t}_{L^2(I_T;H^1)}^2.
 \label{eq:global-base-producer-value}
\end{equation}
The graph identity \eqref{eq:base-graph-identity} on \([0,T]\) gives
\begin{equation}
 \mathfrak B_2[u_0,\nu;T]
 \le4T\norm{u_0}_2^2
 +\kappa_\nu\left(
  \nu Z_0+\mathfrak J_1[u_0,\nu;T]
 \right).
 \label{eq:global-base-producer-graph-bound}
\end{equation}
This value belongs to the unique global history generated by the datum.
Writing it down quantifies that history, but a final estimate must also let a
reader obtain an upper bound without knowing the history in advance.  The
equation supplies that last step.  A proposed history can be inserted into
Navier--Stokes, its defect can be measured, and one scalar comparison then
decides whether the proposal encloses the exact solution on the chosen
horizon.

For each requested finite output, fix once and for all validated rational upper
bounds for every analytic constant used through the requested orders.  We write
these as \(\overline K_q\), \(\overline A_r\),
\(\overline C_{s,k,b}\), \(\overline C_\ell^{\rm Lip}\),
\(\overline G_\ell\), \(\overline C_1^{\rm Lip}\),
\(\overline C_{\rm quad}\), and \(\overline C_{\rm lin}(\nu)\).
Only finitely many are needed for a finite request, and the Fourier,
convolution, embedding, and heat estimates used above give computable rational
upper enclosures for them.  In every expression carrying the superscript
\({\rm alg}\), the corresponding barred constants are used.  For the restart
put
\begin{equation}
 \underline c_\nu
 :=\left(
 8\overline C_{\rm quad}\overline C_{\rm lin}(\nu)^2
 \right)^{-2}
 \le c_\nu.
 \label{eq:algorithmic-restart-constant}
\end{equation}
The unbarred constants remain in the classical estimates.

\begin{proposition}[Residual certificate and terminating datum search]
\label{prop:residual-certificate}
Fix an integer \(s\ge3\).  Enlarge the order constant \(K_s\) from
\eqref{eq:critical-coefficient-row-estimate}, if needed, and fix a proved
computable rational value for it so that, for divergence-free fields \(v,w\),
\begin{align}
 \left|\pair{B(v,w)+B(w,v)}{w}_{H^s}\right|
 &\le K_s\norm v_{H^{s+1}}\norm w_{H^s}^2,
 \notag\\
 \left|\pair{B(w,w)}{w}_{H^s}\right|
 &\le K_s\norm w_{H^s}^3.
 \label{eq:certificate-stability-products}
\end{align}
Let
\[
 v\in C^1([0,T];H_\sigma^s)
       \cap C([0,T];H_\sigma^{s+2})
\]
be any trial history and define its full equation residual by
\begin{equation}
 \operatorname{Res}_\nu[v]
 :=v_t-\nu\Delta v+B(v,v).
 \label{eq:certificate-residual}
\end{equation}
Put
\begin{align}
 A_{s,v}(t)
 &:=K_s\int_0^t\norm{v(\tau)}_{H^{s+1}}\,d\tau,
 \notag\\
 d_{s,v}
 &:=\norm{u_0-v(0)}_{H^s}
 +\int_0^T e^{-A_{s,v}(t)}
   \norm{\operatorname{Res}_\nu[v](t)}_{H^s}\,dt,
 \notag\\
 J_{s,v}(t)
 &:=K_s\int_0^t e^{A_{s,v}(\tau)}\,d\tau,
 \qquad
 \Theta_{s,v}:=1-d_{s,v}J_{s,v}(T).
 \label{eq:certificate-data}
\end{align}
If \(\Theta_{s,v}>0\), then
\begin{equation}
 \norm{u(t)-v(t)}_{H^s}
 \le
 \mathcal E_{s,v}(t)
 :=
 \frac{e^{A_{s,v}(t)}d_{s,v}}
 {1-d_{s,v}J_{s,v}(t)}
 \qquad(0\le t\le T).
 \label{eq:certificate-error-envelope}
\end{equation}

There is a fixed countable family of such divergence-free trial histories
\(\mathscr D_{s,T}=\{v_{s,T}^{(q)}:q\in\N\}\), independent of the solution,
for which the accepted set
\begin{equation}
 \mathcal Q_s[u_0,\nu;T]
 :=\left\{q\in\N:\Theta_{s,v_{s,T}^{(q)}}>0\right\}
 \label{eq:certificate-accepted-set}
\end{equation}
is nonempty.  Its canonical classical selector is
\begin{equation}
 q_s^{\rm cl}[u_0,\nu;T]
 :=\min\mathcal Q_s[u_0,\nu;T].
 \label{eq:certificate-classical-selector}
\end{equation}
This selector and the majorant below are canonical relative to the once-fixed
effective enumeration \(\mathscr D_{s,T}\) and the selected rational value of
\(K_s\).  Those fixed selector choices are suppressed in the notation.
Define
\begin{equation}
 \mathfrak M_s[u_0,\nu;T]
 :=
 \sup_{0\le t\le T}
 \left(
  \norm{v_{s,T}^{(q_s^{\rm cl})}(t)}_{H^s}
  +\mathcal E_{s,v_{s,T}^{(q_s^{\rm cl})}}(t)
 \right).
 \label{eq:certificate-sobolev-majorant}
\end{equation}
Then
\begin{equation}
 \sup_{0\le t\le T}\norm{u(t)}_{H^s}
 \le\mathfrak M_s[u_0,\nu;T]<\infty.
 \label{eq:certificate-sobolev-bound}
\end{equation}
Every quantity on the right of
\eqref{eq:certificate-sobolev-majorant} is evaluated on a finite trial
history and the given datum.  If \(u_0\) has an effective \(H^s\) Cauchy name
and \(\nu>0\) and \(T>0\) are computable reals, a fixed dovetailed validated
search returns an accepted index
\(q_s^{\rm alg}[u_0,\nu;T]\) and a rational upper certificate
\(\overline{\mathfrak M}_s^{\rm alg}[u_0,\nu;T]\) in finite time, with
\begin{equation}
 \sup_{0\le t\le T}\norm{u(t)}_{H^s}
 \le \overline{\mathfrak M}_s^{\rm alg}[u_0,\nu;T].
 \label{eq:certificate-algorithmic-majorant}
\end{equation}
The algorithmic selector need not equal the canonical classical selector.
\end{proposition}

\begin{proof}
Let \(w=u-v\).  Subtracting the residual equation from Navier--Stokes gives
\[
 w_t-\nu\Delta w+B(v,w)+B(w,v)+B(w,w)
 =-\operatorname{Res}_\nu[v].
\]
Pair this equation with \(w\) in \(H^s\).  The viscous term is nonnegative,
and \eqref{eq:certificate-stability-products} gives, after regularizing the
norm at \(w=0\),
\begin{equation}
 D^+\norm w_{H^s}
 \le
 K_s\norm v_{H^{s+1}}\norm w_{H^s}
 +K_s\norm w_{H^s}^2
 +\norm{\operatorname{Res}_\nu[v]}_{H^s}.
 \label{eq:certificate-dini-inequality}
\end{equation}
Set \(y=e^{-A_{s,v}}\norm w_{H^s}\).  Integration of
\eqref{eq:certificate-dini-inequality} gives
\[
 y(t)
 \le d_{s,v}
 +K_s\int_0^t e^{A_{s,v}(\tau)}y(\tau)^2\,d\tau.
\]
Bihari's inequality now yields
\[
 y(t)
 \le
 \frac{d_{s,v}}{1-d_{s,v}J_{s,v}(t)}
\]
as long as the denominator is positive.  The acceptance condition keeps it
positive through \(T\), proving
\eqref{eq:certificate-error-envelope}.

It remains to construct the search family and prove that the search stops.
In Fourier space, fix a uniform effective enumeration of real vector fields
whose transforms are finite rational combinations of computable
\(C_c^\infty\) bumps supported in rational boxes away from the origin.
Apply the solenoidal multiplier
\[
 \widehat{\PP f}(\xi)
 =\left(I-\frac{\xi\otimes\xi}{|\xi|^2}\right)\widehat f(\xi).
\]
The resulting fields are divergence-free and belong to every Sobolev space.
They are dense in \(H_\sigma^m(\R^3)\) for each finite \(m\): smooth compact
Fourier support away from zero is dense in the weighted Fourier \(L^2\)
space, and \(\PP\) is bounded and fixes divergence-free fields.  Rational
time polynomials with coefficients in this spatial family form
\(\mathscr D_{s,T}\).  To see density in the two-norm path space, first apply
smooth spatial Fourier cutoffs to a path and its time derivative; the cutoffs
converge in \(C^1H^s\cap CH^{s+2}\) and place both in every higher Sobolev
space.  Banach-valued polynomial approximation in time, followed by rational
approximation of its finitely many spatial coefficients, then makes
\(\mathscr D_{s,T}\) dense in
\[
 C^1([0,T];H_\sigma^s)
 \cap C([0,T];H_\sigma^{s+2}).
\]

The global solution belongs to this space.  Choose a sequence
\(v_j\in\mathscr D_{s,T}\) converging to \(u\) in its defining norm.  Since
\(\operatorname{Res}_\nu[u]=0\), the product estimate gives
\[
 \norm{\operatorname{Res}_\nu[v_j]}_{C([0,T];H^s)}
 \longrightarrow0,
 \qquad
 \norm{u_0-v_j(0)}_{H^s}\longrightarrow0.
\]
The values \(A_{s,v_j}(T)\), and hence \(J_{s,v_j}(T)\), stay bounded.  Thus
\(d_{s,v_j}J_{s,v_j}(T)\to0\), so some finite candidate is accepted.

For the effective statement, a finite candidate, its nonlinear Fourier
convolution, and its full residual have compact computable Fourier support.
Validated quadrature thus encloses every norm and time integral in
\eqref{eq:certificate-data}; the effective \(H^s\) name of \(u_0\) encloses
the initial mismatch, while the computable \(\nu,T\), and \(K_s\) enclose the
remaining operations.  At stage \(L\), run these enclosures for every
candidate \(q\le L\) to width at most \(2^{-L}\), and accept the first one
whose resulting interval for \(\Theta_{s,v_{s,T}^{(q)}}\) has positive lower
endpoint.  The candidates and enclosure depths are being dovetailed, so no
zero or negative candidate can stall the search.  Some candidate has a
strictly positive margin; at a finite stage its enclosure exposes that margin
and supplies \(q_s^{\rm alg}\).  At that accepting stage let
\[
 \underline\Theta_s>0,
 \qquad
 \overline V_s\ge
 \sup_{0\le t\le T}\norm{v_{s,T}^{(q_s^{\rm alg})}(t)}_{H^s},
 \qquad
 \overline d_s\ge d_{s,v_{s,T}^{(q_s^{\rm alg})}},
 \qquad
 \overline e_s\ge e^{A_{s,v_{s,T}^{(q_s^{\rm alg})}}(T)}
\]
be rational validated bounds from the same enclosure schedule.  Since
\(A_{s,v}\) and \(J_{s,v}\) are increasing,
\[
 \mathcal E_{s,v_{s,T}^{(q_s^{\rm alg})}}(t)
 \le \frac{\overline e_s\,\overline d_s}{\underline\Theta_s}
 \qquad(0\le t\le T).
\]
Thus the search outputs the rational number
\begin{equation}
 \overline{\mathfrak M}_s^{\rm alg}
 :=\overline V_s+
 \frac{\overline e_s\,\overline d_s}{\underline\Theta_s},
 \label{eq:certificate-rational-output}
\end{equation}
which proves \eqref{eq:certificate-algorithmic-majorant}.  This separates the
canonical mathematical data functional from the terminating validated
certificate.
\end{proof}

The order \(s=3\) pays the base rectangle directly.  The projected equation,
\eqref{eq:row-product-constant}, and
\eqref{eq:certificate-sobolev-bound} give
\[
 \norm{u_t(t)}_{H^1}
 \le
 \nu\mathfrak M_3[u_0,\nu;T]
 +A_1\mathfrak M_3[u_0,\nu;T]^2.
\]
Consequently,
\begin{equation}
 \begin{aligned}
 \mathfrak B_2[u_0,\nu;T]
 &\le\mathfrak B_2^{\rm cert}[u_0,\nu;T],\\
 \mathfrak B_2^{\rm cert}[u_0,\nu;T]
 &:=
 T\left\{
 \mathfrak M_3[u_0,\nu;T]^2
 +\left(
 \nu\mathfrak M_3[u_0,\nu;T]
 +A_1\mathfrak M_3[u_0,\nu;T]^2
 \right)^2
 \right\}.
 \end{aligned}
 \label{eq:certified-base-rectangle}
\end{equation}
For effective inputs, the terminating version is obtained from
\eqref{eq:certificate-algorithmic-majorant}:
\begin{equation}
 \mathfrak B_2[u_0,\nu;T]
 \le \overline{\mathfrak B}_2^{\rm alg}[u_0,\nu;T],
 \qquad
 \overline{\mathfrak B}_2^{\rm alg}
 \ge
 T\left\{
  (\overline{\mathfrak M}_3^{\rm alg})^2
  +\left(
   \nu\overline{\mathfrak M}_3^{\rm alg}
   +\overline A_1(\overline{\mathfrak M}_3^{\rm alg})^2
  \right)^2
 \right\},
 \label{eq:algorithmic-base-rectangle}
\end{equation}
where the right member is a validated rational upper enclosure of the
displayed computable-real expression.  The canonical value
\(\mathfrak B_2^{\rm cert}\) is a datum-to-number map for every Schwartz
datum; \(\overline{\mathfrak B}_2^{\rm alg}\) is the independently checkable
finite output for effective data.  Both constructions are downstream from
global continuation, whose existence guarantees acceptance on every finite
 horizon.  Every monotone certified recursion below also has its terminating
 algorithmic version after replacing \(\mathfrak B_2^{\rm cert}\) by
 \(\overline{\mathfrak B}_2^{\rm alg}\) and every analytic constant it uses
 by the corresponding validated rational upper bound fixed above.  Its restart
 output uses \(\underline c_\nu\).

On the same horizon define
\begin{equation}
 \mathfrak C_{1/2}[u_0,\nu;T]
 :=\sup_{0\le S\le T}\int_0^S\mathcal P_{1/2}(t)\,dt,
 \qquad
 \mathfrak A_{1/2}[u_0,\nu;T]
 :=\frac{H(0)+\mathfrak C_{1/2}[u_0,\nu;T]}{\nu}.
 \label{eq:global-critical-current-values}
\end{equation}
The exact critical balance gives
\begin{equation}
 \int_0^T\norm{\Lambda^{3/2}u(t)}_2^2\,dt
 \le\mathfrak A_{1/2}[u_0,\nu;T].
 \label{eq:global-critical-coefficient-value}
\end{equation}
The base rectangle gives the sharper history-resolved coefficient
\begin{equation}
 \mathfrak A_{1/2}^{\rm prod}[u_0,\nu;T]
 :=\min\left\{
 \mathfrak B_2[u_0,\nu;T],
 \left(\frac{\norm{u_0}_2^2}{2\nu}\right)^{3/4}
 \mathfrak B_2[u_0,\nu;T]^{1/4}
 \right\},
 \label{eq:global-producer-critical-coefficient}
\end{equation}
while the certificate supplies the data-resolved coefficient
\begin{equation}
 \mathfrak A_{1/2}^{\rm cert}[u_0,\nu;T]
 :=\min\left\{
 \mathfrak B_2^{\rm cert}[u_0,\nu;T],
 \left(\frac{\norm{u_0}_2^2}{2\nu}\right)^{3/4}
 \mathfrak B_2^{\rm cert}[u_0,\nu;T]^{1/4}
 \right\}.
 \label{eq:global-certified-critical-coefficient}
\end{equation}
Monotonicity and \eqref{eq:certified-base-rectangle} give
\begin{equation}
 \int_0^T\norm{\Lambda^{3/2}u(t)}_2^2\,dt
 \le\mathfrak A_{1/2}^{\rm prod}[u_0,\nu;T]
 \le\mathfrak A_{1/2}^{\rm cert}[u_0,\nu;T].
 \label{eq:global-producer-critical-coefficient-bound}
\end{equation}

Repeat the finite definitions
\eqref{eq:spatial-envelope-exponent}--\eqref{eq:explicit-rectangle-envelope}
with \(T\) in place of \(T_*\).  Using
\(\mathfrak A_{1/2}^{\rm prod}\) gives the unadorned quantities
\[
 \Gamma_m(T),\quad
 \mathfrak S_m(T),\quad
 \mathfrak D_m(T),\quad
 \mathfrak V_{n,m}(T),\quad
 \mathfrak Q_{n,m}(T),\quad
 \mathfrak R_{N,M}(T).
\]
Using \(\mathfrak A_{1/2}^{\rm cert}\) gives
\[
 \Gamma_m^{\rm cert}(T),\quad
 \mathfrak S_m^{\rm cert}(T),\quad
 \mathfrak D_m^{\rm cert}(T),\quad
 \mathfrak V_{n,m}^{\rm cert}(T),\quad
 \mathfrak Q_{n,m}^{\rm cert}(T),\quad
 \mathfrak R_{N,M}^{\rm cert}(T).
\]
Every operation in these definitions is nonnegative and increasing, so each
unadorned value is bounded by its certified counterpart.

The \(r=1\) product row also retains the two analytic feedback bounds.  Put
\begin{align}
 \mathfrak J_1^\sharp(T)
 &:=
 \frac{G_1}{\sqrt2\,\nu}
 \sqrt{\norm{u_0}_2^2\,D_1(u_0)}\,e^{\Gamma_1(T)/2}
 \left(
  D_1(u_0)e^{\Gamma_1(T)}
  +D_2(u_0)e^{\Gamma_2(T)}
 \right),
 \notag\\
 \mathfrak J_1^{\mathrm{fin}}(T)
 &:=
 \min\left\{
  A_1^2T\mathfrak S_2(T)^4,\,
  \mathfrak J_1^\sharp(T)
 \right\}.
 \label{eq:global-graph-action-best}
\end{align}
Then
\begin{equation}
 \begin{aligned}
 \mathfrak B_2[u_0,\nu;T]
 &\le\overline{\mathfrak B}_2[u_0,\nu;T],\\
 \overline{\mathfrak B}_2[u_0,\nu;T]
 &:=4T\norm{u_0}_2^2
 +\kappa_\nu\left(
  \nu Z_0+\mathfrak J_1^{\mathrm{fin}}(T)
 \right).
 \end{aligned}
 \label{eq:global-closed-base-rectangle}
\end{equation}
This feedback form can be sharper once the history has been evaluated.
Equation~\eqref{eq:certified-base-rectangle} is the unconditional
datum-resolved majorant used in the final statement.

Whenever one finite order \(\ell\ge3\) satisfies
\(\Xi_\ell[u_0,\nu;T]<1\),
Proposition~\ref{prop:finite-order-datum-window} also gives the closed
short-horizon values
\begin{equation}
 \mathfrak C_{1/2}[u_0,\nu;T]
 \le\widehat{\mathfrak C}_{1/2}^{(\ell)}[u_0,\nu;T],
 \qquad
 \mathfrak A_{1/2}[u_0,\nu;T]
 \le\widehat{\mathfrak A}_{1/2}^{(\ell)}[u_0,\nu;T].
 \label{eq:global-finite-order-data-majorants}
\end{equation}
Repeating the same recursions with
\(\widehat{\mathfrak A}_{1/2}^{(\ell)}\) gives hatted quantities.  On that
window the strongest displayed estimate is the minimum of the certified and
hatted majorants.

\begin{corollary}[Complete quantitative output]
\label{cor:complete-quantitative-output}
For every \(N,M,n,m,k\in\Nzero\), every \(1\le a\le\infty\), and every
\(2\le b\le\infty\), the global solution satisfies the following
datum-resolved estimates on \([0,T]\).

Every finite rectangle and every uniform jet row obey
\begin{align}
 \mathcal N_{N,M}(u;I_T)
 &\le\mathfrak R_{N,M}(T)
 \le\mathfrak R_{N,M}^{\rm cert}(T),
 \label{eq:final-rectangle-bound}\\
 \sup_{0\le t\le T}\norm{\partial_t^nu(t)}_{H^m}
 &\le\mathfrak V_{n,m}(T)
 \le\mathfrak V_{n,m}^{\rm cert}(T),
 \label{eq:final-velocity-row-bound}\\
 \sup_{0\le t\le T}\norm{\partial_t^np(t)}_{H^m}
 &\le\mathfrak Q_{n,m}(T)
 \le\mathfrak Q_{n,m}^{\rm cert}(T).
 \label{eq:final-pressure-row-bound}
\end{align}
If \(s\in\Nzero\) satisfies
\(s>k+\frac32-\frac3b\), then
\begin{align}
 \norm{\nabla^k\partial_t^nu}_{L^a(I_T;L^b)}
 &\le C_{s,k,b}T^{1/a}\mathfrak V_{n,s}^{\rm cert}(T),
 \label{eq:final-mixed-jet-bound}\\
 \norm{\nabla^k\partial_t^np}_{L^a(I_T;L^b)}
 &\le C_{s,k,b}T^{1/a}\mathfrak Q_{n,s}^{\rm cert}(T),
 \label{eq:final-pressure-jet-bound}
\end{align}
where \(T^{1/\infty}=1\).  The same right sides without the factor
\(T^{1/a}\) control the corresponding \(L_t^\infty L_x^b\) norms.

For \(j\in\Nzero\), \(m>j+\frac12\), and \(K\ge1\), define
\begin{equation}
 E_{j+1/2,>K}(t)
 :=\frac12\int_{|\xi|>K}|\xi|^{2j+1}
 |\widehat u(t,\xi)|^2\,d\xi.
 \label{eq:final-critical-tail-definition}
\end{equation}
The complete critical-height ladder and every one of its frequency tails
satisfy
\begin{align}
 \sup_{0\le t\le T}E_{j+1/2}(t)
 &\le\frac12\mathfrak S_{j+1}^{\rm cert}(T)^2,
 \label{eq:final-critical-height-sup}\\
 \int_0^T E_{j+1/2}(t)\,dt
 &\le\frac12\mathfrak D_j^{\rm cert}(T)^2,
 \label{eq:final-critical-height-mass}\\
 \sup_{0\le t\le T}E_{j+1/2,>K}(t)
 &\le\frac12K^{2j+1-2m}
 D_m(u_0)e^{\Gamma_m^{\rm cert}(T)}.
 \label{eq:final-frequency-tail}
\end{align}
For the original critical row,
\begin{align}
 \norm{H'}_{L^1(0,T)}
 &\le\frac12\mathfrak B_2^{\rm cert}[u_0,\nu;T],
 \notag\\
 \norm{\mathcal P_{1/2}}_{L^1(0,T)}
 &\le\left(\frac12+\nu\right)
       \mathfrak B_2^{\rm cert}[u_0,\nu;T].
 \label{eq:final-critical-action}
\end{align}

The endpoint at the chosen horizon is quantitative.  For \(0\le t\le T\),
\begin{align}
 \norm{\partial_t^nu(t)-\partial_t^nu(T)}_{H^m}
 &\le(T-t)\mathfrak V_{n+1,m}^{\rm cert}(T),
 \label{eq:final-endpoint-modulus}\\
 \norm{\partial_t^np(t)-\partial_t^np(T)}_{H^m}
 &\le(T-t)(\mathfrak L^Q_{n,m})^{\rm cert}(T),
 \label{eq:final-pressure-endpoint-modulus}
\end{align}
where \((\mathfrak L^Q_{n,m})^{\rm cert}(T)\) is
\eqref{eq:pressure-endpoint-modulus-constant} evaluated with the certified
row values.  Restarting the solution from \(u(T)\) has the explicit lower
lifespan
\begin{equation}
 \delta_T
 \ge\min\left\{1,
 c_\nu\bigl(1+\mathfrak S_3^{\rm cert}(T)\bigr)^{-2}
 \right\}.
 \label{eq:final-restart-lifespan}
\end{equation}

The derivative ledger remains finite and exact.  The certificate itself reads
the datum in \(H^3\).  The velocity row \((n,m)\) additionally uses datum
order \(m+2n\), while the pressure row uses order
\(\sigma_m+2n\).  The rectangle \((N,M)\) uses order \(M+2N+1\);
retaining \(Q_N\) raises this, when needed, to
\[
 \max\{M+2N+1,\sigma_M+2N\}.
\]
The velocity endpoint modulus uses order \(m+2n+2\), and its pressure
counterpart uses order \(\sigma_m+2n+2\).  The uniform bound for
\(E_{j+1/2}\) uses datum order \(j+1\), while its time integral uses order
\(j\).  The order-\(m\) frequency tail and the restart use datum orders
\(m\) and \(3\).  In particular, the displayed mixed velocity estimate at
embedding order \(s\) reads depth \(\max\{3,s+2n\}\), while its pressure
 counterpart reads depth \(\max\{3,\sigma_s+2n\}\).  Thus every certified
 datum depth is
\(\max\{3,d\}\), with \(d\) the order just listed.  The hatted route reads
depth \(\max\{\ell,d\}\).  Numerical evaluation requires an effective Sobolev
name through that finite depth together with computable names for \(\nu\) and
\(T\).  It is also relative to the fixed trial-history enumeration, selector
constant \(K_s\), and the validated rational analytic constants specified
above.  Temporal order, spatial order, viscosity, horizon, and datum depth all
remain visible.

If \(\Xi_\ell[u_0,\nu;T]<1\), the recursive jet, height, tail, and endpoint
upper bounds may be replaced by the smaller of the certified value and its
hatted counterpart.  In particular,
\begin{equation}
 \mathcal N_{N,M}(u;I_T)
 \le\min\left\{
 \mathfrak R_{N,M}^{\rm cert}(T),
 \widehat{\mathfrak R}_{N,M}^{(\ell)}(T)
 \right\}.
 \label{eq:final-data-only-window}
\end{equation}
The critical-action pair uses
\begin{equation}
 \mathfrak B_{2,\ell}^{\rm best}(T)
 :=\min\left\{
  \mathfrak B_2^{\rm cert}[u_0,\nu;T],
  \widehat{\mathfrak R}_{0,2}^{(\ell)}(T)
 \right\},
 \label{eq:final-best-critical-action-coefficient}
\end{equation}
so that
\begin{equation}
 \norm{H'}_{L^1(0,T)}
 \le\frac12\mathfrak B_{2,\ell}^{\rm best}(T),
 \qquad
 \norm{\mathcal P_{1/2}}_{L^1(0,T)}
 \le\left(\frac12+\nu\right)\mathfrak B_{2,\ell}^{\rm best}(T).
 \label{eq:final-best-critical-action}
\end{equation}
The restart estimate is a lower bound, so its two versions combine in the
opposite direction:
\begin{align}
 \delta_T
 &\ge\max\left\{
  \min\left\{1,c_\nu(1+\mathfrak S_3^{\rm cert}(T))^{-2}\right\},
  \min\left\{1,c_\nu(1+\widehat{\mathfrak S}_3^{(\ell)}(T))^{-2}\right\}
 \right\}
 \notag\\
 &=\min\left\{1,
  c_\nu\left(1+\min\left\{
   \mathfrak S_3^{\rm cert}(T),
   \widehat{\mathfrak S}_3^{(\ell)}(T)
  \right\}\right)^{-2}
 \right\}.
 \label{eq:final-best-restart-lifespan}
\end{align}
For effective inputs, replacing \(c_\nu\) by \(\underline c_\nu\) and the
upper envelopes by their validated rational enclosures gives a terminating
rational lower bound for the same restart time.
\end{corollary}

\begin{proof}
Equation~\eqref{eq:certified-base-rectangle} and monotonicity of
\(x\mapsto\min\{x,cx^{1/4}\}\) give
\eqref{eq:global-producer-critical-coefficient-bound}.  Substituting the
certified coefficient into
\eqref{eq:critical-coefficient-row-estimate}--\eqref{eq:closed-rectangle-estimate}
gives \eqref{eq:final-rectangle-bound}--\eqref{eq:final-pressure-row-bound}.
Sobolev embedding gives
\eqref{eq:final-mixed-jet-bound}--\eqref{eq:final-pressure-jet-bound}.
The comparisons
\[
 |\xi|^{2j+1}\le(1+|\xi|^2)^{j+1},
 \qquad
 |\xi|^{2j+1}\mathbf1_{\{|\xi|>K\}}
 \le K^{2j+1-2m}|\xi|^{2m},
\]
together with the certified spatial envelopes, give
\eqref{eq:final-critical-height-sup}--\eqref{eq:final-frequency-tail}.
The variation and response estimates give
\eqref{eq:final-critical-action}.  Integrating the next velocity and pressure
rows gives the two endpoint moduli, and the local construction in
Appendix~\ref{app:local-theory} gives
\eqref{eq:final-restart-lifespan}.  Under
\(\Xi_\ell[u_0,\nu;T]<1\), the alternative coefficient
\eqref{eq:global-finite-order-data-majorants} enters the same nonnegative
recursions.  Taking the smaller upper output proves
\eqref{eq:final-data-only-window} and
\eqref{eq:final-best-critical-action}; monotonic decrease of the restart map
in its \(H^3\) input proves \eqref{eq:final-best-restart-lifespan}.
\end{proof}
